\documentclass[11pt,a4paper]{article}
\usepackage[utf8]{inputenc}
\usepackage[T1]{fontenc}
\usepackage[spanish,es-noshorthands,es-nodecimaldot,es-tabla]{babel}
\usepackage{lmodern}
\usepackage[margin=2.2cm]{geometry}
\usepackage{amsmath,amssymb}
\usepackage{graphicx}
\usepackage{booktabs}
\usepackage{float}
\usepackage[font=small,labelfont=bf]{caption}
\usepackage{xcolor}
\usepackage[colorlinks=true,linkcolor=blue!50!black,citecolor=blue!50!black,urlcolor=blue!50!black]{hyperref}
\usepackage{enumitem}
\setlist{nosep,leftmargin=*}

\graphicspath{{resultados/figs/}}
\newcommand{\fig}[1]{Fig.~\ref{#1}}
\newcommand{\ok}{\textcolor{green!50!black}{\textbf{sí}}}
\newcommand{\no}{\textcolor{red!70!black}{\textbf{no}}}
% Resultado pendiente: recuadro gris en lugar de la figura o del valor
\newcommand{\pend}{\textcolor{gray}{\textit{pendiente}}}
\newcommand{\figpend}[2]{%
  \IfFileExists{resultados/figs/#1.pdf}{\includegraphics[width=#2]{#1.pdf}}{%
    \fbox{\parbox[c][4cm][c]{\dimexpr#2-2\fboxsep-2\fboxrule}{\centering\color{gray}\large Simulación en curso:\\ figura \texttt{\detokenize{#1}} pendiente}}}}

\title{Sincronismo BiSAR por camino directo sin referencia externa\\[2pt]
\large Vuelo circular de los dos radares: análisis de viabilidad y validación por simulación}
\author{David Atencia}
\date{31 de octubre de 2026}

\begin{document}
\maketitle

\begin{abstract}
\noindent
Dos drones vuelan el mismo círculo de 50~m de radio a 100~m de altura, a la
misma velocidad y separados 90$^\circ$ en azimut. El Radar~1 transmite y
recibe (monostático); el Radar~2 es un receptor pasivo sin GPS, PPS ni reloj
común que graba de forma discontinua el camino directo y los ecos. Se
estudia si el camino directo basta para sincronizarlo con la precisión que
exige el Backprojection (BP) bistático sobre la vuelta completa, y se valida
con una simulación MATLAB de extremo a extremo. \textbf{Resultado:} sí es
viable si se usa la \emph{fase portadora} del camino directo en cada ventana
grabada (estrategia E3). Con E3 la imagen es indistinguible de la de un
reloj perfecto (pérdida $<0.02$~dB, desplazamiento $<1$~mm, $\rho=1.000$),
con cualquier oscilador (de OCXO a un cristal barato) y con cualquier patrón
de escucha, incluso grabando solo el 1\,\% de las PRT. El sincronismo solo
de tiempo pierde $\approx3$~dB, y un modelo de reloj solo funciona con OCXO.
Requisitos: SNR del directo $\gtrsim20$~dB y por encima de los ecos;
\emph{stagger} de emisión $\ge1~\mu$s registrado en el log del Tx (la
distancia Tx--Rx2 es constante, así que la geometría no ayuda a identificar
los pulsos; no hace falta pulso de inicio); al menos 8 ventanas con directo;
navegación de cada dron con error $\lesssim3$~cm, que es menos exigente que
lo que pide el monostático.

\smallskip\noindent
Las estrategias comparadas se llaman E0 a E3 (explicación sencilla en la
sección~\ref{sec:alg}): E0 pone el reloj en hora una sola vez, E1 predice
su desviación con un modelo, E2 mide en cada pulso el instante de llegada
del directo y E3 mide en cada pulso el instante y la fase del directo.
\end{abstract}

\tableofcontents
\newpage

\begin{figure}[H]
\centering
\includegraphics[width=\textwidth]{inf_conclusiones.pdf}
\caption{Resumen gráfico de los resultados.
(a)~Error de fase residual frente a la SNR del directo: teoría (CRB, líneas)
y simulación (marcadores). La banda verde es la región en la que la imagen
pierde menos de 1~dB ($\sigma_\phi<0.47$~rad). E3 entra en ella desde
20~dB (por debajo de 15~dB el directo no se distingue de los ecos) y E2
solo por encima de $\sim$35~dB.
(b)~Pérdida de pico según el oscilador del Rx2: E3 no pierde nada con
ninguno, E2 pierde $\approx3$~dB con todos y E1 solo funciona con OCXO.
(c)~Probabilidad de asociar bien los pulsos según el \emph{stagger} $S$ y el
número de ventanas $J$ (TCXO): con PRF fija es imposible.
(d)~Patrones de escucha: pérdida de pico debida al sincronismo (barras) y
cambio de PSLR por muestreo de la apertura (línea negra). E3 no pierde nada
en ningún patrón; E2 pierde $\approx3$~dB en todos.
(e)~Error al interpolar la fase en huecos sin camino directo (simulado y
teoría de puente browniano).
(f)~Pérdida de pico frente al error de navegación de ambos drones: E3
tolera $\approx3$~cm, más que el monostático.}
\label{fig:conclusiones}
\end{figure}
\newpage

%=====================================================================
\section{Escenario simulado}
\label{sec:escenario}

Los dos drones vuelan el \emph{mismo} círculo (centro en el origen, radio
$R=50$~m, altura $h=100$~m) a 10~m/s, con el Radar~2 adelantado 90$^\circ$
en azimut (\fig{fig:escenario}a,b). Como giran con la misma velocidad
angular, la distancia entre ellos es constante:
$R_d=2R\sin(45^\circ)=70.7$~m. El Rx2 escucha una vuelta completa (31.4~s,
12\,566 PRT) y la escena tiene cinco reflectores puntuales en el suelo,
dentro de la proyección del círculo.

Con una apertura circular completa el soporte espectral sobre el suelo es un
anillo (\fig{fig:escenario}c) y la respuesta al impulso es de tipo $J_0$
(\fig{fig:escenario}d): lóbulo principal de 0.30~m a $-3$~dB, primer nulo a
0.43~m y anillos que decaen despacio (PSLR $-7.8$~dB). Esto es intrínseco a
la geometría circular, no al sincronismo. Por eso las métricas de imagen se
comparan siempre con la imagen formada con el reloj ideal.

\begin{figure}[H]
\centering
\includegraphics[width=0.95\textwidth]{inf_escenario.pdf}
\caption{Geometría de la simulación. (a)~Vista 3D: círculo de vuelo común,
posición de los radares en $t_0$ (Tx círculo azul, Rx2 triángulo naranja),
camino directo (discontinuo) y camino bistático al blanco central (verde).
(b)~Planta con las curvas de igual exceso de camino bistático $R_b-R_d$ en
$t_0$ y las posiciones en $t_0+T/8$ y $t_0+T/4$ (atenuadas). La línea
Tx--Rx2 gira con los drones pero conserva su longitud. (c)~Soporte espectral
del blanco central a lo largo de la vuelta: un anillo de radio
$k_r=5.63$~rad/m. (d)~Respuesta al impulso BiSAR con reloj ideal.}
\label{fig:escenario}
\end{figure}

\begin{table}[H]
\centering\footnotesize
\caption{Parámetros de la simulación.}
\label{tab:param}
\begin{tabular}{@{}ll@{\qquad}ll@{}}
\toprule
Portadora $f_c$ & 425~MHz ($\lambda=0.705$~m) & Vuelo & círculo común, $R=50$~m, $h=100$~m \\
Ancho de banda $B$ & 50~MHz & Velocidad & 10~m/s (ambos), 31.4~s por vuelta \\
Duración $T_p$ & 1.3~$\mu$s & Separación & Rx2 adelantado 90$^\circ$, $R_d=70.7$~m \\
PRF & 400~Hz, stagger $\pm5~\mu$s & Escucha & 12\,566 PRT (una vuelta) \\
Muestreo & 100~MHz complejo & SNR directo / eco & 30 / 20~dB por pulso \\
Oscilador por defecto & TCXO (1~ppm) & Offset inicial $\delta_0$ & 5--20~s (arbitrario) \\
Resolución & 0.30~m ($-3$~dB), nulo 0.43~m & PSLR / ISLR ideal & $-7.8$ / $+7.8$~dB (IRF $J_0$) \\
\bottomrule
\end{tabular}
\end{table}

La \fig{fig:senal} muestra lo que graba el Rx2. Cada ventana contiene el
camino directo, siempre el primero en llegar y unos 10~dB más fuerte que los
ecos, que llegan entre 140 y 190~m de exceso de camino. Ese pico es la
referencia de tiempo y de fase de toda la ventana. El Rx2 no sabe qué pulso
está escuchando: solo tiene marcas de su propio reloj, que en una vuelta
deriva unos 20~$\mu$s con un TCXO y además está mezclado con el
\emph{stagger} de la emisión (\fig{fig:senal}d, puntos azules). Tras
asociar las ventanas con el log del Tx, la deriva aparece limpia (línea
naranja).

\begin{figure}[H]
\centering
\includegraphics[width=\textwidth]{inf_senal.pdf}
\caption{Señal en el Rx2. (a)~Patrones de escucha discontinua simulados (cada
columna es una PRT; en color, las PRT grabadas). (b)~Salida del filtro
adaptado de una ventana: directo a 30~dB y ecos de los cinco blancos a 20~dB.
(c)~Dato comprimido de las 12\,566 ventanas, alineado con el directo y con
el directo cancelado: las pistas de los ecos oscilan con la vuelta.
(d)~Llegadas del directo medidas con el reloj local, quitando la PRT
nominal: el \emph{stagger} ($\pm5~\mu$s) oculta la deriva. Tras la
asociación queda el error de reloj $\delta(t)$ (TCXO,
$\varepsilon=0.64$~ppm, 20~$\mu$s en la vuelta).}
\label{fig:senal}
\end{figure}

%=====================================================================
\section{Modelo de reloj y parámetros a estimar}
\label{sec:modelo}

El reloj del Radar~1 es la referencia $t$. El Rx2 lleva un tiempo local
\begin{equation}
t_2=t+\delta(t),\qquad
\delta(t)=\delta_0+\varepsilon t+\tfrac12\gamma t^2+w(t)+j_k ,
\end{equation}
con offset inicial $\delta_0$, error de frecuencia $\varepsilon$, deriva
$\gamma$, ruido del oscilador $w(t)$ (paseo aleatorio,
$\operatorname{Var}[w(t+\tau)-w(t)]=\sigma_y^2(1\,\mathrm{s})\,\tau$) y jitter
$j_k$ por ventana. Si el LO y el reloj de muestreo salen del mismo oscilador,
la fase de demodulación es $\varphi_2=2\pi f_c\,\delta(t)+\phi_0$. El offset
de tiempo, el de frecuencia ($\Delta f=\varepsilon f_c$), la deriva y el
offset de fase son un único proceso $\delta(t)$, visto desde la envolvente y
desde la portadora.

Para cada pulso usado, el BP necesita la hora local $t_n+\tau_b+\delta$ y la
fase $2\pi f_c\delta+\phi_0$. Hay que estimar el índice de pulso de cada
ventana (asociación), $\delta_0,\varepsilon,\gamma$ y, sobre todo, la parte
aleatoria $w$ en cada ventana, que es la que fija la coherencia de fase. Con
31~s de apertura esa parte aleatoria pesa más que en una apertura corta
(\fig{fig:reloj}).

\begin{figure}[H]
\centering
\includegraphics[width=\textwidth]{inf_reloj.pdf}
\caption{Por qué no basta un modelo de reloj a lo largo de una vuelta.
Arriba: error de reloj $\delta(t)-\delta_0$ para las cuatro clases de
oscilador (de 0.9 a 940~$\mu$s en 31~s). Abajo: error de fase que queda tras
ajustar perfectamente $\delta_0+\varepsilon t+\gamma t^2/2$ (lo mejor que
puede hacer E1); la banda verde es la tolerancia de 0.47~rad. Solo el OCXO
queda dentro (0.03~rad). El TCXO ya la cruza (0.34~rad rms, picos de
0.8~rad), y XO y XO barato quedan muy fuera.}
\label{fig:reloj}
\end{figure}

%=====================================================================
\section{Observabilidad}
\label{sec:obs}

En cada ventana $j$ el directo da el retardo y la fase
\begin{equation}
\hat a_j=t_{n_j}+\tau_d(n_j)+\delta(t_j)+\nu_j,\qquad
\hat\varphi_j=-2\pi f_c[\tau_d(n_j)+\delta(t_j)]-\phi_0+\eta_j .
\end{equation}
De ahí salen cinco condiciones:
\begin{enumerate}
\item \textbf{Distancia Tx--Rx2 conocida.} $\tau_d$ y $\delta$ aparecen
  sumados: un error en la distancia entre drones no se distingue de un error
  de reloj. Para la fase debe conocerse con error $\ll\lambda$, a partir de
  la navegación de ambos drones.
\item \textbf{Índice de pulso.} Desplazar la hipótesis $k$ pulsos cambia las
  llegadas en $k\,\mathrm{PRT}+(s_{n+k}-s_n)+[\tau_d(n+k)-\tau_d(n)]$. En el
  vuelo circular común $\tau_d$ es \emph{constante}, así que la geometría no
  aporta ninguna firma: con PRF uniforme $k$ es inobservable. El
  \emph{stagger} pseudoaleatorio $s_n$ registrado por el Tx es
  imprescindible.
\item \textbf{Modelo de reloj.} $\varepsilon$ necesita al menos 2 ventanas
  separadas y $\gamma$ al menos 3;
  $\sigma_\varepsilon\approx\sigma_t\sqrt{12/J}/T_{obs}$.
\item \textbf{Fase por ventana.} $2\pi f_c\delta(t_j)+\phi_0$ es observable
  módulo $2\pi$ en cada ventana con directo, que es justo lo que necesita el
  BP. No hace falta desenrollar la fase entre ventanas.
\item \textbf{Aproximaciones.} $\varepsilon f_cT_p\ll1$ y
  $\varepsilon T_{\text{ventana}}\ll1/B$, ciertas incluso con 20~ppm. El
  Doppler del directo es nulo, porque la distancia Tx--Rx2 no cambia.
\end{enumerate}

%=====================================================================
\section{Precisión necesaria para un BiSAR enfocado}
\label{sec:tol}

Se inyectan errores controlados en la base de tiempos y se mide la imagen
(\fig{fig:tol}). En apertura circular completa la \emph{posición} del
blanco la fija la fase a lo largo de la vuelta, no la envolvente. Por eso
las leyes cambian respecto a una apertura rectilínea:
\begin{itemize}
\item un offset de tiempo constante \textbf{no desplaza} la imagen: solo
  hace muestrear el pulso comprimido fuera de su pico, con una tolerancia
  amplia (7.9~ns);
\item un error de fase de baja frecuencia (frecuencia residual o deriva)
  se reparte en una componente sinusoidal con el azimut, que el BP convierte
  en \emph{desplazamiento}, y un residuo que \emph{desenfoca}. Se predicen
  ambos ajustando el error de fase con
  $c_0+k\,\nabla R_b(\theta)\cdot\Delta\mathbf p$ a lo largo de la vuelta
  (líneas discontinuas de la \fig{fig:tol}b).
\end{itemize}
Como $f_c/B=8.5$ y el pico lo fija la fase, \textbf{el requisito que domina
es el de fase}.

\begin{figure}[H]
\centering
\includegraphics[width=\textwidth]{exp1_tolerancias.pdf}
\caption{Tolerancias (exp1). (a)~Offset de tiempo: solo pérdida de
envolvente, sin desplazamiento. (b)~Frecuencia residual: pérdida (eje
izquierdo) y desplazamiento (eje derecho) simulados frente a la teoría de
primer orden. (c)~Deriva cuadrática: la pérdida es pequeña (1~dB con
5.6~rad), pero desplaza el pico 0.05~m ya con 0.61~rad; lo que limita es el
desplazamiento. (d)~Ruido de fase
blanco: $e^{-\sigma^2/2}$. (e)~Jitter de tiempo: equivale a una fase
$2\pi f_c\sigma_t$. Líneas rojas: criterio de $-1$~dB.}
\label{fig:tol}
\end{figure}

\begin{table}[H]
\centering\small
\caption{Tolerancias para 425~MHz/50~MHz en vuelo circular (criterio:
desplazamiento $\le0.05$~m, $\approx0.1$ de la resolución, o pérdida
$\le1$~dB).}
\label{tab:tol}
\begin{tabular}{@{}lll@{}}
\toprule
Error & Efecto & Tolerancia \\
\midrule
Offset de tiempo constante & pérdida de envolvente (sin desplazamiento) & 7.9~ns \\
Frecuencia residual $\Delta\varepsilon$ & desplazamiento + desenfoque & $1.0\cdot10^{-11}$ (0.004~Hz) \\
Fase cuadrática (borde) & desplazamiento (+ desenfoque) & 0.61~rad (desplazamiento) \\
Fase blanca & pérdida de pico & 0.47~rad rms \\
Jitter de tiempo & ídem vía $2\pi f_c\sigma_t$ & \textbf{176~ps rms} \\
Índice de pulso erróneo & imagen destruida & debe ser 0 \\
\bottomrule
\end{tabular}
\end{table}

La envolvente es muy fácil de cumplir, pero la fase exige \textbf{176~ps
de coherencia pulso a pulso} y una frecuencia residual de $10^{-11}$, cinco
veces más estricta que con una apertura rectilínea de 6~s, porque la vuelta
dura 31~s. Un TCXO libre (1~ppm) está $10^5$ veces por encima.

%=====================================================================
\section{Cómo explotar el camino directo}
\label{sec:alg}

\begin{enumerate}
\item \textbf{Adquisición:} la primera ventana graba una PRT completa y se
  detecta el \emph{primer} pico que supera 13~dB, que es el directo.
\item \textbf{Puertas de seguimiento:} cada ventana se centra en la última
  llegada medida del directo más $m\,\mathrm{PRT}$. El margen cubre el
  stagger y la deriva máxima del reloj durante el hueco desde esa medida
  ($\pm13~\mu$s con escucha continua).
\item \textbf{Estimación ML del directo} (Newton sobre el filtro adaptado
  interpolado). Alcanza las cotas de Cramér--Rao
  \begin{equation}
  \sigma_t=\frac{1}{2\pi\beta\sqrt{2\,\mathrm{SNR}}},\qquad
  \sigma_\phi=\frac{1}{\sqrt{2\,\mathrm{SNR}}},\qquad \beta=12.7~\text{MHz}.
  \end{equation}
\item \textbf{Asociación:} para cada hipótesis $k$ se ajusta por mínimos
  cuadrados un modelo de reloj sobre las llegadas predichas con el log y se
  elige la de menor residuo.
\item \textbf{Base de tiempos y fase por ventana:}
  $\hat\delta_j=\hat a_j-t_{n_j}-\tau_d(n_j)$ y
  $\hat\psi_j=-[\arg\hat r_j+2\pi f_c\tau_d(n_j)]$.
\item \textbf{Cancelación del directo y BP} sobre la vuelta completa (sin
  ponderación en azimut, porque la apertura circular cerrada no tiene bordes):
  $I(\mathbf g)=\sum_j\mathrm{RC}_j(t_{n_j}+\tau_b+\hat\delta_j)\,e^{j(2\pi f_c\tau_b+\hat\psi_j)}$.
\end{enumerate}

\begin{table}[H]
\centering\small
\caption{Estrategias comparadas (ninguna usa verdad terreno).}
\label{tab:estr}
\begin{tabular}{@{}llll@{}}
\toprule
 & Retardo & Fase & Error de fase esperado \\
\midrule
E0 & 1.ª ventana, constante & constante & $2\pi f_c\varepsilon t$ (falla) \\
E1 & modelo $\delta_0+\varepsilon t+\gamma t^2/2$ & $2\pi f_c\hat\delta(t)$ & ruido de oscilador no modelado \\
E2 & medido en cada ventana & $2\pi f_c\hat\delta_j$ & $(f_c/\beta)/\sqrt{2\,\mathrm{SNR}}$ \\
E3 & medido en cada ventana & \textbf{fase portadora medida} & $1/\sqrt{2\,\mathrm{SNR}}$ \\
\bottomrule
\end{tabular}
\end{table}

E3 aplica el principio del canal de referencia del radar pasivo. Es mejor
que E2 en un factor $(f_c/\beta)^2=1124$, es decir, \textbf{30.5~dB de SNR}.

\subsection*{Las cuatro estrategias en palabras simples}

El Rx2 tiene un reloj que no coincide con el del Tx: empezó en otro
momento, va un poco más rápido o más lento y además tiembla de forma
aleatoria. Para formar la imagen hay que saber, en cada pulso, \emph{cuánto
se ha desviado ese reloj}. La única pista disponible es el camino directo:
el pulso que llega del Tx al Rx2 sin rebotar en el suelo, cuya distancia sí
conocemos. Las cuatro estrategias son cuatro formas cada vez más cuidadosas
de usar esa pista.

\begin{description}[leftmargin=1.2em, style=nextline]
\item[E0 --- ``poner el reloj en hora una vez'' (offset)]
  Se mide el desfase con el primer pulso directo y se supone que ya no
  cambia. Es como poner en hora un reloj de pulsera barato al salir de casa y
  fiarse de él todo el día: al poco tiempo va adelantado o atrasado. En 31~s
  el reloj del Rx2 se desvía microsegundos, y la imagen \textbf{no se forma}.
  Sirve como referencia de lo que pasa si no se sincroniza.

\item[E1 --- ``aprender cómo se desvía el reloj'' (modelo)]
  Se usan todos los pulsos directos para ajustar una curva suave del error
  del reloj: un desfase inicial, más un adelanto constante (el reloj va
  rápido o lento), más una deriva lenta. Es como saber que un reloj se
  adelanta 2~s al día y corregirlo con esa regla. Funciona si el reloj es
  muy estable (OCXO). Con relojes peores falla, porque además de la parte
  predecible el reloj \emph{tiembla} de forma aleatoria y la curva suave no
  puede seguir esos temblores.

\item[E2 --- ``mirar la hora en cada pulso'' (solo tiempo)]
  En cada ventana grabada se mide de nuevo \emph{cuándo} llega el pulso
  directo y se corrige el reloj con esa medida. Ya no se predice nada: se
  mide siempre, así que sigue también los temblores. El problema es la
  precisión de la medida. El instante de llegada se mide con un error de
  unos 300~ps, y para enfocar a 425~MHz hace falta mejor que 176~ps. Es como
  mirar la hora en un reloj que solo marca minutos cuando necesitas
  segundos: la imagen se forma, pero pierde unos 3~dB de brillo.

\item[E3 --- ``mirar la hora y el segundero en cada pulso'' (tiempo + fase)]
  Además de cuándo llega el pulso directo, se mide su \emph{fase}: en qué
  punto de su oscilación llega la onda portadora de 425~MHz, que vibra cada
  2.35~ns. Esa medida es unas 30 veces más fina que la del instante de
  llegada (equivale a $\sim$10~ps frente a 300~ps), porque se mira la onda en sí y no solo
  la forma del pulso. Y como el pulso directo y los ecos de la escena llegan
  a la misma ventana, a través del mismo receptor y del mismo reloj, la fase
  del directo sirve directamente de referencia para los ecos: cualquier
  error del reloj en esa ventana afecta igual a los dos y se cancela. Es la
  técnica del canal de referencia de los radares pasivos. Con E3 la imagen
  es \textbf{igual a la de un reloj perfecto}.
\end{description}

En resumen: E0 corrige una vez, E1 predice, E2 mide el instante en cada
pulso y E3 mide el instante y la fase en cada pulso. Solo E3 es lo bastante
precisa para cualquier oscilador; E1 solo sirve con osciladores muy
estables.

%=====================================================================
\section{Demostración de extremo a extremo}
\label{sec:e2e}

Una realización completa con TCXO ($\varepsilon=0.64$~ppm), un offset
inicial arbitrario, escucha aleatoria del 50\,\% de las PRT (6251 de
12\,566) y SNR del directo de 30~dB. La asociación encuentra el pulso
correcto con un margen de $2\cdot10^3$ entre la mejor hipótesis y la segunda.

\begin{figure}[H]
\centering
\includegraphics[width=\textwidth]{exp2_reloj_asociacion.pdf}
\caption{Sincronismo en la realización de demostración. Arriba: error de
reloj verdadero y error de la base de tiempos estimada. Abajo izquierda:
error de fase por ventana: E2 sale de la banda de tolerancia, E1 deriva
lentamente con el paseo aleatorio del oscilador y E3 es casi nulo. Abajo
derecha: función de coste de la asociación, con un único mínimo.}
\label{fig:e2e}
\end{figure}

\begin{figure}[H]
\centering
\includegraphics[width=\textwidth]{exp2_imagenes.pdf}
\caption{Imágenes BP de la escena (escala común, 35~dB). E0 no forma imagen.
E1, E2 y E3 enfocan los cinco blancos con la IRF anular de la apertura
circular. E2 pierde 2.8~dB de pico, que en la escena se ve como un fondo más
alto. Abajo: imagen monostática y fusión coherente mono~+~E3. La diferencia
de fase entre los picos mono y bi es de 0$^\circ$: E3 ancla también la fase
absoluta.}
\label{fig:imgs}
\end{figure}

\begin{figure}[H]
\centering
\includegraphics[width=\textwidth]{inf_cortes.pdf}
\caption{Cortes de la respuesta del blanco central. E3 (morado) se superpone
exactamente al reloj ideal (negro discontinuo). E2 baja 3~dB todo el corte.
E1 apenas pierde pico, pero cambia ligeramente los anillos. E0 queda por
debajo de $-45$~dB y no aparece.}
\label{fig:cortes}
\end{figure}

\begin{table}[H]
\centering\small
\caption{Métricas de la demostración, comparadas con la imagen de reloj
ideal (PSLR e ISLR medidos en un parche de $\pm3$~m; la IRF $J_0$ tiene
anillos altos incluso con reloj ideal).}
\label{tab:e2e}
\begin{tabular}{@{}lrrrrrrrc@{}}
\toprule
 & $\sigma_t$ (ps) & $\sigma_\phi$ (rad) & despl. (m) & pérdida (dB) & PSLR & ISLR & $\rho$ & ¿Enfoca? \\
\midrule
E0 & $1.2\cdot10^7$ & 1.81 & 2.67 & $-63.0$ & $-0.1$ & 12.8 & 0.19 & \no \\
E1 & 79 & 0.21 & 0.017 & $-0.20$ & $-8.8$ & 3.4 & 0.982 & \ok \\
E2 & 298 & 0.80 & 0.005 & $-2.76$ & $-8.5$ & 6.0 & 0.997 & \no \\
E3 & 298 & \textbf{0.024} & 0.000 & $\mathbf{-0.00}$ & $-7.7$ & 7.8 & 1.000 & \ok \\
Ideal & 0 & 0 & 0 & 0 & $-7.7$ & 7.8 & 1 & -- \\
\bottomrule
\end{tabular}
\end{table}

%=====================================================================
\section{Precisión alcanzada: SNR del directo y oscilador}
\label{sec:snr}

Monte Carlo de la cadena completa sobre la vuelta (4 realizaciones por
caso, escucha continua). Primero se barre la SNR del directo con TCXO
(\fig{fig:snr}) y después se cambia el oscilador a SNR$_d$~=~30~dB
(\fig{fig:osc}). La asociación fue correcta en todas las realizaciones.

\begin{figure}[H]
\centering
\figpend{exp3_snr}{\textwidth}
\caption{Barrido de SNR del directo (TCXO). Izquierda: el error de retardo
alcanza la CRB desde 20~dB, pero solo baja de la tolerancia de 176~ps por
encima de $\sim$35~dB. Centro: el error de fase de E3 sigue la CRB de fase,
$\sim$30~dB por debajo del de E2. Derecha: pérdida de pico de la imagen. E3
cumple desde 15~dB ($-0.9$~dB) y es prácticamente perfecta desde 20~dB; E2
necesita $\ge40$~dB. Por debajo de 15~dB el detector toma un eco por el
directo y todo falla.}
\label{fig:snr}
\end{figure}

\begin{table}[H]
\centering\small
\caption{Barrido de SNR del directo (TCXO, medianas de 4 realizaciones).
$\sigma_t$: error del retardo medido; $\sigma_\phi$: error de la corrección
de fase. Última columna: fracción de realizaciones que cumplen todos los
criterios.}
\label{tab:snr}
\begin{tabular}{@{}rrrrrrrrc@{}}
\toprule
SNR$_d$ (dB) & $\sigma_t$ (ps) & CRB$_t$ (ps) & $\sigma_\phi$ E2 & $\sigma_\phi$ E3 & pérd. E1 & pérd. E2 & pérd. E3 & cumple E1/E2/E3 \\
\midrule
10 & $5\cdot10^5$ & 3001 & 1.81 & 1.47 & $-47.9$ & $-42.6$ & $-29.4$ & 0 / 0 / 0 \\
15 & $2\cdot10^5$ & 1688 & 1.80 & 0.65 & $-29.7$ & $-32.5$ & $-0.92$ & 0 / 0 / \textbf{1} \\
20 & 955 & 949 & 1.77 & 0.076 & $-0.30$ & $-28.3$ & $-0.04$ & 0.75 / 0 / \textbf{1} \\
30 & 301 & 300 & 0.80 & 0.024 & $-0.39$ & $-2.82$ & $-0.00$ & 1 / 0 / \textbf{1} \\
40 & 95 & 95 & 0.25 & 0.008 & $-0.42$ & $-0.29$ & $-0.00$ & 0.5 / 1 / \textbf{1} \\
50 & 30 & 30 & 0.08 & 0.002 & $-0.53$ & $-0.03$ & $-0.00$ & 0.5 / 1 / \textbf{1} \\
\bottomrule
\end{tabular}
\end{table}

\begin{figure}[H]
\centering
\figpend{exp3_oscilador}{0.8\textwidth}
\caption{Clase de oscilador a SNR$_d$~=~30~dB. Izquierda: error de fase.
E3 da 0.024~rad con los cuatro osciladores: lo fija la SNR, no el reloj.
E2 da 0.8~rad con todos: lo fija la precisión del retardo. E1 empeora con
el ruido del oscilador. Derecha: pérdida de pico. E1 cumple con OCXO
($-0.01$~dB), está en el límite con TCXO ($-0.45$~dB, cumple en la mitad
de las realizaciones) y falla con XO ($-11$~dB) y XO barato ($-29$~dB).}
\label{fig:osc}
\end{figure}

Lectura de los resultados:
\begin{itemize}
\item El estimador del directo \textbf{alcanza la cota teórica} desde 20~dB,
  en tiempo y en fase.
\item \textbf{Por debajo de 15~dB lo que falla es la detección, no la
  estimación}: los ecos (20~dB por pulso en esta escena) superan al directo
  y el detector toma un eco por el directo. En el diseño real el directo
  tiene que ser claramente más fuerte que los ecos, lo habitual con línea de
  vista entre drones a 70~m.
\item \textbf{E3 es independiente del oscilador}; E1 solo es fiable con OCXO.
  Con TCXO, E1 pierde poco pico pero falla a veces por lóbulos, porque el
  error de fase del paseo aleatorio acumulado en 31~s es de baja frecuencia.
\end{itemize}

%=====================================================================
\section{Pulsos perdidos (escucha discontinua)}
\label{sec:huecos}

Hay que separar tres efectos. (i)~\textbf{Sobre el sincronismo E3, ninguno}:
cada ventana grabada se referencia con su propio directo. (ii)~Sobre la
imagen hay un efecto de muestreo de la apertura, que no es un problema de
sincronismo. Con PRF de 400~Hz la vuelta está muy sobremuestreada en azimut
($\Delta\theta=5\cdot10^{-4}$~rad por PRT), y los anillos espurios de un
patrón periódico 1/$G$ aparecen a $\sim2\pi/(k_r\,G\,\Delta\theta)$ del
blanco (17~m con $G=128$): fuera del parche de medida, pero visibles en la
escena completa. (iii)~Si una ventana no contiene el directo, su fase hay
que interpolarla, y el error crece como un puente browniano:
\begin{equation}
\sigma_\phi(T_g)\approx\sqrt{(2\pi f_c\sigma_y)^2T_g/6+\sigma_w^2},\qquad
T_{g,\max}\approx 6\left(\frac{0.47}{2\pi f_c\sigma_y}\right)^2 .
\end{equation}

\begin{figure}[H]
\centering
\includegraphics[width=\textwidth]{exp4_pulsos_perdidos.pdf}
\caption{Pulsos perdidos (TCXO, 30~dB, 3 realizaciones). Izquierda: pérdida
de pico frente a la fracción escuchada. E3 queda en 0~dB (peor caso
$-0.02$~dB), E1 en $-0.3$~dB y E2 en $\approx-3$~dB, su límite de SNR.
Centro: los lóbulos dentro del parche apenas cambian por el muestreo
periódico ($<0.5$~dB); con escucha muy escasa y aleatoria (1\,\%) o en
ráfagas cortas el PSLR sube 2--3~dB, por muestreo y no por sincronismo. Derecha: error de fase interpolado en huecos sin directo,
frente a la teoría. Hueco máximo: OCXO/TCXO $>1.3$~s (límite del ensayo),
XO 0.16~s, XO barato menos de 2 PRT.}
\label{fig:huecos}
\end{figure}

%=====================================================================
\section{Asociación de pulsos}
\label{sec:asoc}

\subsection*{Qué es el stagger, con un ejemplo}

El \emph{stagger} es un pequeño retardo, distinto en cada pulso, que el Tx
añade a su instante de emisión y apunta en su log. Con PRF de 400~Hz los
pulsos deberían salir cada 2.5~ms exactos; con stagger de $\pm5~\mu$s salen,
por ejemplo, en $0+1.3~\mu$s, $2.5~\text{ms}-3.7~\mu$s,
$5.0~\text{ms}+0.4~\mu$s, $7.5~\text{ms}+4.1~\mu$s\ldots{} Sin stagger
todos los pulsos son iguales y equiespaciados, y el Rx2 no puede saber si
escuchó los pulsos 0, 1, 2 o los 517, 518, 519. Con stagger, las
separaciones entre llegadas que mide el Rx2 (por ejemplo $+4.1$, $+3.7$,
$-6.3~\mu$s respecto a 2.5~ms) forman una huella que solo coincide con un
tramo del log. Esa huella identifica los pulsos y, de paso, el desfase
inicial del reloj: \textbf{no hace falta un pulso de inicio}. Son $\pm5~\mu$s
sobre 2.5~ms (0.2\,\%), y el BP usa el instante real de cada emisión, así
que no afecta a la imagen.

\subsection*{Resultados}

Modelo ligero (sin formar señales): llegadas del directo con el reloj
completo del Rx2 (offset, frecuencia, deriva y paseo aleatorio) y un error de
medida igual a la CRB a 30~dB (300~ps), con $J$ ventanas repartidas al azar
en la vuelta. Para cada hipótesis de pulso la asociación ajusta un modelo de
reloj \emph{cuadrático} (desfase + frecuencia + deriva). Con 31~s de vuelo
el ajuste lineal de la versión anterior no bastaba: la deriva no modelada
del reloj actuaba como ruido y con XO barato solo acertaba el 73\,\% de las
veces con stagger de 5~$\mu$s y 5 ventanas (98\,\% con el cuadrático).

\begin{figure}[H]
\centering
\figpend{exp5_asociacion}{0.85\textwidth}
\caption{Probabilidad de asociación correcta frente al número de ventanas
para varios \emph{stagger} (200 ensayos por punto), con TCXO y con XO
barato. Sin stagger es imposible con cualquier número de ventanas.}
\label{fig:asoc}
\end{figure}

\begin{table}[H]
\centering\small
\caption{Probabilidad de asociación correcta (TCXO / XO barato) según el
stagger $S$ y el número de ventanas $J$.}
\label{tab:asoc}
\begin{tabular}{@{}lccccc@{}}
\toprule
Stagger $S$ & $J=3$ & $J=4$ & $J=5$ & $J=8$ & $J=16$ \\
\midrule
0 (PRF fija) & 0 / 0 & 0 / 0 & 0 / 0 & 0 / 0 & 0 / 0 \\
100~ns & 0.04 / 0 & 0.13 / 0.01 & 0.97 / 0.04 & 1 / 0.80 & 1 / 1 \\
1~$\mu$s & 0.16 / 0.01 & 0.53 / 0.07 & 1 / 0.74 & 1 / 1 & 1 / 1 \\
5~$\mu$s & 0.46 / 0.02 & 0.88 / 0.21 & 1 / 0.98 & 1 / 1 & 1 / 1 \\
\bottomrule
\end{tabular}
\end{table}

Con $S\ge1~\mu$s y al menos \textbf{8 ventanas} con directo, la asociación
es correcta en el 100\,\% de los casos con cualquier oscilador. En la
práctica el Rx2 graba cientos o miles de ventanas por vuelta, así que el
margen es enorme; el requisito real es que el Tx emita con stagger y lo
registre.

%=====================================================================
\section{Error de navegación de los drones}
\label{sec:nav}

Los dos drones vuelan con un error de navegación aleatorio (AR(1),
$\tau_c=1$~s), independiente en cada uno y de la misma $\sigma$ por eje. Se
compara cada imagen con la misma adquisición sin error de navegación
(\fig{fig:nav}).

El camino directo mide reloj y distancia Tx--Rx2 a la vez, así que E3
toma la parte del error de navegación proyectada sobre la línea Tx--Rx2 como
si fuera error de reloj y la ``corrige'' en todos los ecos de esa ventana.
En la geometría rectilínea anterior (solo el Tx con error) esto compensaba
parte del error. Aquí ocurre lo contrario: la línea Tx--Rx2 es horizontal,
los ecos llegan casi en vertical, y el error que se resta no es el que
afecta a los ecos. Por eso E3 tolera algo \emph{menos} que el bistático con
reloj perfecto.

\begin{figure}[H]
\centering
\figpend{exp6_navegacion}{0.8\textwidth}
\caption{Pérdida de pico (izquierda) y coherencia con la imagen sin error
(derecha) frente al error de navegación de cada dron (3 realizaciones).
Error tolerable para perder 1~dB: $\approx3.3$~cm con E3, $\approx4$~cm
con reloj ideal y $\approx2.8$~cm en el monostático. E2 parte de su pérdida
propia de 2.8~dB.}
\label{fig:nav}
\end{figure}

\begin{table}[H]
\centering\small
\caption{Pérdida de pico (dB, peor blanco) frente al error de navegación de
cada dron, comparada con la misma adquisición sin error.}
\label{tab:nav}
\begin{tabular}{@{}rrrrr@{}}
\toprule
$\sigma_{nav}$ por eje & Bi E2 & \textbf{Bi E3} & Bi reloj ideal & Mono \\
\midrule
1~cm & $-2.89$ & $-0.09$ & $-0.07$ & $-0.14$ \\
2~cm & $-3.21$ & $-0.35$ & $-0.27$ & $-0.58$ \\
5~cm & $-4.95$ & $\mathbf{-2.08}$ & $-1.49$ & $-2.99$ \\
10~cm & $-9.99$ & $-7.11$ & $-5.67$ & $-9.43$ \\
\bottomrule
\end{tabular}
\end{table}

En la práctica la navegación no es un obstáculo adicional para el
sincronismo. Con E3 se toleran unos 3~cm por dron, algo más que lo que ya
exige el monostático (2.8~cm). El requisito de navegación lo sigue fijando
el monostático, no la sincronización.

%=====================================================================
\section{Conclusiones de viabilidad}
\label{sec:concl}

\begin{enumerate}
\item \textbf{Es viable con la vuelta completa.} Usando solo el camino
  directo, la estrategia E3 produce imágenes BiSAR indistinguibles de las
  del reloj ideal (pérdida $<0.02$~dB, desplazamiento $<1$~mm, $\rho=1.000$;
  \fig{fig:imgs}, \fig{fig:cortes}) con arranque y offset arbitrarios y con
  cualquier patrón de escucha, hasta el 1\,\% de las PRT
  (\fig{fig:huecos}).
\item \textbf{El sincronismo solo de tiempo no basta.} El directo da
  $\sim$300~ps a 30~dB y la tolerancia es de 176~ps: E2 pierde
  $\approx3$~dB. La fase portadora es $(f_c/\beta)^2\approx30$~dB más
  precisa.
\item \textbf{Condiciones de SNR:} el directo debe llegar con
  SNR~$\gtrsim20$~dB por pulso y claramente por encima de los ecos. E3
  alcanza la cota teórica desde 20~dB; por debajo de 15~dB falla la
  detección (\fig{fig:snr}).
\item \textbf{El modelo de reloj (E1) solo sirve con OCXO.} En 31~s el ruido
  del oscilador acumula demasiada fase: con TCXO E1 está en el límite
  ($-0.45$~dB, cumple la mitad de las veces) y con XO o XO barato falla por
  completo. E3 da el mismo resultado con los cuatro osciladores
  (\fig{fig:osc}).
\item \textbf{La geometría circular no ayuda a la asociación, pero el stagger
  sí.} La distancia Tx--Rx2 es constante: sin \emph{stagger} el índice de
  pulso es inobservable. Con stagger $\ge1~\mu$s registrado en el log del Tx
  y al menos 8 ventanas con directo, la asociación es correcta siempre, con
  cualquier oscilador. No hace falta pulso de inicio (\fig{fig:asoc}).
\item \textbf{La envolvente es irrelevante:} un offset de tiempo de hasta
  7.9~ns cuesta menos de 1~dB y no desplaza la imagen. Lo que manda es la
  coherencia de fase: 176~ps rms pulso a pulso y $10^{-11}$ de frecuencia
  residual.
\item \textbf{Navegación:} con error en los dos drones, E3 tolera unos 3~cm
  por eje y dron para perder 1~dB, algo menos que el bistático con reloj
  perfecto (4~cm) pero más que el monostático (2.8~cm). El requisito de
  navegación lo sigue fijando el monostático (\fig{fig:nav}).
\end{enumerate}

%=====================================================================
\section{Métricas para demostrar el sincronismo experimentalmente}
\label{sec:exp}

En un ensayo real no se dispone de un reloj ideal con el que comparar.
\begin{itemize}
\item \textbf{Sincronismo, sin verdad terreno:} margen de asociación ($\gg10$);
  residuo del retardo del directo tras el modelo frente a la CRB a la SNR
  medida; validación cruzada (ajustar con las ventanas impares y predecir las
  pares); desviación de Allan de $\hat\delta(t)$ frente a la hoja de datos;
  consistencia de $\hat\varepsilon$ entre vueltas; dispersión residual de la
  fase del directo. En vuelo circular, además, la distancia Tx--Rx2 medida
  por el directo debe ser constante.
\item \textbf{Imagen, con reflectores triédricos de posición conocida:} error
  de posición de los picos ($<0.05$~m); forma anular de la IRF frente a la
  teórica de la geometría; coherencia entre semivueltas; el autofocus no
  debería mejorar la imagen; diferencia de fase mono--bi en los picos y
  ganancia de la fusión coherente; repetibilidad entre vueltas.
\item \textbf{Criterios usados aquí} (peor blanco, frente al reloj ideal):
  pérdida $\le1$~dB, desplazamiento $\le0.05$~m, $\Delta$PSLR y
  $\Delta$ISLR $\le3$~dB, $\rho\ge0.95$.
\end{itemize}

%=====================================================================
\section{Limitaciones del modelo}
\label{sec:lim}

\begin{itemize}
\item Blancos puntuales isótropos, sin diagramas de antena ni clutter. En
  vuelo circular real la antena del Rx2 tendría que ver a la vez al Tx y a la
  escena.
\item Sin multicamino en el directo. Con ambos drones a 100~m, el rebote en
  el suelo del directo tiene 141~m de exceso de camino: se separa bien del
  directo ($\gg c/B$), pero cae dentro de la zona de los ecos de la escena.
\item Sin saturación ni cuantización del ADC (directo fuerte y ecos débiles).
\item La navegación del Rx2 se registra con su propio reloj y se lleva al
  tiempo del Tx con la base de tiempos sincronizada; con $\delta$ conocido a
  ns, el error de posición resultante es despreciable.
\item Oscilador con FM blanca y deriva. El ruido flicker empeoraría E1, pero
  no afecta a E3.
\end{itemize}

\vspace{1em}
\noindent\small\sloppy Reproducción: \texttt{run\_all.m} (experimentos) y
\texttt{fig\_informe.m} (figuras adicionales de este informe) en
\texttt{simulations/analise\_sincronismo\_3110/}. Los datos de cada figura
están en los CSV de \texttt{resultados/}; los resultados de la geometría
rectilínea anterior se conservan en \texttt{resultados/antiguo\_lineal/}.

\end{document}
